\exercisesheader{}

% 1

\eoce{\qt{Luas di bawah kurva, Bagian I\label{area_under_curve_1}} Berapa persen
dari distribusi normal baku $N(\mu=0, \sigma=1)$ yang terdapat pada setiap
daerah berikut? Pastikan Anda menggambar grafik. \vspace{-3mm}
\begin{multicols}{4}
\begin{parts}
\item $Z < -1.35$
\item $Z > 1.48$
\item $-0.4 < Z < 1.5$
\item $|Z| > 2$
\end{parts}
\end{multicols}
}{}

% 2

\eoce{\qt{Luas di bawah kurva, Bagian II\label{area_under_curve_2}} Berapa persen
dari distribusi normal baku $N(\mu=0, \sigma=1)$ yang terdapat pada setiap
daerah berikut? Pastikan Anda menggambar grafik. \vspace{-3mm}
\begin{multicols}{4}
\begin{parts}
\item $Z > -1.13$
\item $Z < 0.18$
\item $Z > 8$
\item $|Z| < 0.5$
\end{parts}
\end{multicols}
}{}

% 3

\eoce{\qt{Skor GRE, Bagian I\label{GRE_intro}} Sophia mengikuti Graduate Record
Examination (GRE) dan memperoleh skor 160 pada bagian Penalaran Verbal
(Verbal Reasoning, VR) serta 157 pada bagian Penalaran Kuantitatif
(Quantitative Reasoning, QR). Rata-rata skor bagian Penalaran Verbal untuk
seluruh peserta ujian adalah 151 dengan simpangan baku 7, sedangkan rata-rata
skor Penalaran Kuantitatif adalah 153 dengan simpangan baku 7.67.
Anggap kedua distribusi hampir normal.
\begin{parts}
\item Tuliskan notasi ringkas untuk kedua distribusi normal ini.
\item Berapa skor-Z Sophia pada bagian Penalaran Verbal? Pada bagian
Penalaran Kuantitatif? Gambarlah kurva distribusi normal baku dan tandai
kedua skor-Z tersebut.
\item Apa yang ditunjukkan oleh kedua skor-Z ini?
\item Dibandingkan dengan peserta lain, pada bagian mana ia berprestasi
lebih baik?
\item Carilah persentil skornya untuk kedua bagian ujian.
\item Berapa persen peserta ujian yang memperoleh skor lebih tinggi
daripada Sophia pada bagian Penalaran Verbal? Pada bagian Penalaran
Kuantitatif?
\item Jelaskan mengapa hanya membandingkan skor mentah dari kedua bagian
dapat menghasilkan kesimpulan yang keliru tentang bagian mana yang
dikerjakan seorang peserta dengan lebih baik.
\item Jika distribusi skor pada kedua bagian ujian tidak hampir normal,
apakah jawaban Anda untuk bagian (b)--(f) akan berubah? Jelaskan alasan Anda.
\end{parts}
}{}

% 4

\eoce{\qt{Waktu triatlon, Bagian I\label{triathlon_times_intro}} Dalam triatlon,
peserta biasanya ditempatkan dalam kelompok usia dan gender. Leo dan Mary,
yang berteman, sama-sama menyelesaikan Hermosa Beach Triathlon. Leo
bertanding dalam kelompok \textit{Laki-laki, Usia 30--34}, sedangkan Mary
bertanding dalam kelompok \textit{Perempuan, Usia 25--29}. Leo menyelesaikan
perlombaan dalam 1:22:28 (4948 detik), sedangkan Mary menyelesaikannya dalam
1:31:53 (5513 detik). Leo jelas finis lebih cepat, tetapi mereka ingin tahu
bagaimana hasil mereka jika dibandingkan dengan kelompok masing-masing.
Dapatkah Anda membantu mereka? Berikut beberapa informasi tentang kinerja
kedua kelompok:
\begin{itemize}
\setlength{\itemsep}{0mm}
\item Waktu finis kelompok \textit{Laki-laki, Usia 30--34} memiliki rata-rata
4313 detik dengan simpangan baku 583 detik.
\item Waktu finis kelompok \textit{Perempuan, Usia 25--29} memiliki rata-rata
5261 detik dengan simpangan baku 807 detik.
\item Distribusi waktu finis kedua kelompok mendekati distribusi normal.
\end{itemize}
Ingat: kinerja yang lebih baik bersesuaian dengan waktu finis yang lebih cepat.
\begin{parts}
\item Tuliskan notasi ringkas untuk kedua distribusi normal ini.
\item Berapa skor-Z waktu finis Leo dan Mary? Apa yang ditunjukkan
oleh skor-Z tersebut?
\item Siapa yang memiliki peringkat lebih baik dalam kelompoknya,
Leo atau Mary? Jelaskan alasan Anda.
\item Berapa persen peserta triatlon dalam kelompok Leo yang finis
lebih lambat daripada Leo?
\item Berapa persen peserta triatlon dalam kelompok Mary yang finis
lebih lambat daripada Mary?
\item Jika distribusi waktu finis tidak hampir normal, apakah jawaban Anda
untuk bagian (b)--(e) akan berubah? Jelaskan alasan Anda.
\end{parts}
}{}

% 5

\eoce{\qt{Skor GRE, Bagian II\label{GRE_cutoffs}} Pada
Latihan~\ref{GRE_intro}, kita melihat dua distribusi skor GRE:
$N(\mu=151, \sigma=7)$ untuk bagian verbal dan
$N(\mu=153, \sigma=7.67)$ untuk bagian kuantitatif. Gunakan informasi ini
untuk menghitung setiap nilai berikut:
\begin{parts}
\item Skor seorang peserta yang berada pada persentil ke-$80$
di bagian Penalaran Kuantitatif.
\item Skor seorang peserta yang memperoleh hasil lebih buruk daripada
70\% peserta ujian di bagian Penalaran Verbal.
\end{parts}
}{}

\D{\newpage}

% 6

\eoce{\qt{Waktu triatlon, Bagian II\label{triathlon_times_cutoffs}} Pada
Latihan~\ref{triathlon_times_intro}, kita melihat dua distribusi waktu
triatlon: $N(\mu=4313, \sigma=583)$ untuk
\emph{Laki-laki, Usia 30--34} dan $N(\mu=5261, \sigma=807)$ untuk kelompok
\emph{Perempuan, Usia 25--29}. Waktu dinyatakan dalam detik.
Gunakan informasi ini untuk menghitung setiap nilai berikut:
\begin{parts}
\item Batas waktu bagi 5\% atlet tercepat dalam kelompok laki-laki,
yakni atlet dengan 5\% waktu penyelesaian tersingkat.
\item Batas waktu bagi 10\% atlet paling lambat dalam kelompok perempuan.
\end{parts}
}{}

% 7

\eoce{\qt{Cuaca LA, Bagian I\label{la_weather_intro}} Rata-rata suhu maksimum
harian pada bulan Juni di LA adalah 77\degree F dengan simpangan baku
5\degree F. Anggap suhu pada bulan Juni mendekati distribusi normal.
\begin{parts}
\item Berapa peluang mengamati suhu 83\degree F atau lebih tinggi di LA
pada suatu hari yang dipilih secara acak pada bulan Juni?
\item Seberapa rendah suhu pada 10\% hari terdingin (hari dengan suhu
maksimum terendah) selama bulan Juni di LA?
\end{parts}
}{}

% 8

\eoce{\qt{CAPM\label{CAPM}} Capital Asset Pricing Model (CAPM) adalah model
keuangan yang mengasumsikan imbal hasil suatu portofolio berdistribusi normal.
Misalkan suatu portofolio memiliki imbal hasil tahunan rata-rata 14.7\%
(yakni keuntungan rata-rata 14.7\%) dengan simpangan baku 33\%.
Imbal hasil 0\% berarti nilai portofolio tidak berubah, imbal hasil negatif
berarti portofolio merugi, dan imbal hasil positif berarti portofolio
memperoleh keuntungan.
\begin{parts}
\item Dalam berapa persen tahun portofolio ini merugi, yakni memiliki
imbal hasil kurang dari 0\%?
\item Berapa batas untuk 15\% imbal hasil tahunan tertinggi pada
portofolio ini?
\end{parts}
}{}

% 9

\eoce{\qt{Cuaca LA, Bagian II\label{la_weather_unit_change}}
Latihan~\ref{la_weather_intro} menyatakan bahwa rata-rata suhu maksimum
harian pada bulan Juni di LA adalah 77\degree F dengan simpangan baku
5\degree F, dan suhu tersebut dapat diasumsikan mengikuti distribusi normal.
Kita menggunakan persamaan berikut untuk mengubah \degree F (Fahrenheit)
menjadi \degree C (Celsius):
\[ C = (F - 32) \times \frac{5}{9}. \]
\begin{parts}
\item Tuliskan model peluang untuk distribusi suhu dalam \degree C
pada bulan Juni di LA.
\item Berapa peluang mengamati suhu 28\degree C (yang kira-kira
bersesuaian dengan 83\degree F) atau lebih tinggi pada bulan Juni di LA?
Hitung menggunakan model \degree C dari bagian (a).
\item Apakah Anda memperoleh jawaban yang sama atau berbeda pada bagian (b)
soal ini dan bagian (a) Latihan~\ref{la_weather_intro}? Apakah Anda terkejut?
Jelaskan.
\item Perkirakan IQR suhu (dalam \degree C) pada bulan Juni di LA.
\end{parts}
}{}

% 10

\eoce{\qt{Cari simpangan baku\label{find_sd_cholesterol}}
Kadar kolesterol perempuan berusia 20 sampai 34 tahun mengikuti
distribusi yang mendekati normal dengan rata-rata 185 miligram
per desiliter (mg/dl).
Perempuan dengan kadar kolesterol di atas 220 mg/dl dianggap memiliki
kolesterol tinggi, dan sekitar 18.5\% perempuan termasuk dalam kategori ini.
Berapa simpangan baku distribusi kadar kolesterol perempuan
berusia 20 sampai~34 tahun?
}{}
